\section{Weisfeiler--Leman refinement and the bijective game}
\label{sec:wl}

We first fix the refinement convention and its game characterization.
The game will let us query the constraints encoded in the graphs of the
next section. It also gives the monotonicity needed to interpret a
dimension bound as identification by a particular refinement algorithm.

All graphs in this paper are finite, simple, undirected, and uncolored
unless colors are explicitly introduced during a construction. For a
graph $G$, an integer $k\ge 2$, and a tuple
$\bar v=(v_1,\ldots,v_k)\in V(G)^k$, let
$\bar v[i\leftarrow z]$ denote the tuple obtained by replacing its
$i$th entry by $z$. The initial color $\chi_0^G(\bar v)$ records every
equality $v_i=v_j$ and adjacency $v_iv_j\in E(G)$. Define subsequent
colors by
\begin{equation}
\label{eq:wl-update}
\chi_{s+1}^G(\bar v)=
\left(\chi_s^G(\bar v),
\left\{\!\left\{
  \bigl(\chi_s^G(\bar v[1\leftarrow z]),\ldots,
        \chi_s^G(\bar v[k\leftarrow z])\bigr)
       :z\in V(G)
\right\}\!\right\}\right).
\end{equation}
Thus a single multiset records the \emph{joint vector} of replacement
colors. It is not a list of separately counted coordinate multisets.
For $k=1$, we instead use ordinary color refinement: all vertices have
one initial color, and
\[
\chi_{s+1}^G(v)=
\left(\chi_s^G(v),
\left\{\!\left\{\chi_s^G(w):w\in N_G(v)\right\}\!\right\}\right).
\]
The colors are named consistently between graphs. Write $G\equiv_k H$
if, at every round, the multisets of colors of their $k$-tuples agree.
We say that $k$-WL \emph{identifies} $G$ if $G\equiv_k H$ implies
$G\cong H$ for every graph $H$. The Weisfeiler--Leman dimension of $G$
is the least positive $k$ that identifies it.

\subsection{The game characterization}

For graphs of equal positive order, the \emph{bijective game with $K$
pair slots} is played as follows. A position consists of at most $K$
pairs of marked vertices, one vertex from each graph in each occupied
slot. Repeated vertices are allowed. Spoiler may discard pairs. To
fill a vacant slot, Spoiler first designates that slot; Duplicator
then announces a bijection $f:V(G)\to V(H)$; finally, Spoiler chooses
$v\in V(G)$ and places the pair $(v,f(v))$ in the designated slot.
Duplicator wins if the occupied pairs always define a partial
isomorphism, preserving equality as well as adjacency. A winning
strategy must work against every sequence of Spoiler's choices.

This is a standard characterization of the joint-vector convention;
see, for example, \cite{KN2019}. We include the argument to fix both
the parameter shift and the treatment of vacant slots.

\begin{lemma}
\label{lem:wl-game}
Let $G,H$ have equal positive order, and let $k\ge 2$. Then
$G\equiv_k H$ if and only if Duplicator has a winning strategy from
the empty position in the bijective game with $k+1$ pair slots.
Moreover, a winning strategy with at least two pair slots implies
$G\equiv_1 H$.
\end{lemma}

\begin{proof}
Apply refinement with common color names to the disjoint union of the
two \emph{tuple tables} $V(G)^k$ and $V(H)^k$, taking replacements
within the corresponding graph. Since each update retains the old
color, these common partitions refine. If one update causes no
splitting, it only renames the existing classes. All subsequent
updates then have the same property, so the common partition has
stabilized.

We use two consequences of the update rule. First, if two tuples have
equal colors at a round, applying the same coordinate permutation to
both preserves color equality. This follows by induction from the
initial types: a common permutation permutes the components of a
replacement vector and applies the corresponding coordinate
permutation to each replacement tuple. Second, for two tuples in the
same stable class, stability provides a vertex bijection matching
their entire vectors of stable replacement colors. Indeed, the two
multisets of such vectors agree, so we may match the vertices within
each vector class.

Suppose $G\equiv_k H$. Choose a pair of equally colored stable
$k$-tuples; one exists because the color histograms agree and the
graphs are nonempty. Use the preceding bijection to append one
vertex to each tuple. The resulting pair of $(k+1)$-tuples has the
following property: deleting any common slot leaves equally colored
stable $k$-tuples. Deleting the appended slot gives the original
pair; every other deletion gives a matched replacement pair, up to
a common coordinate permutation.

Duplicator maintains a pair of fully specified $(k+1)$-tuples with
this property, treating entries in vacant slots as virtual. To fill
slot $j$, apply the stable-class bijection to the projection deleting
slot $j$. Whichever vertex Spoiler chooses, the projection deleting
$j$ remains unchanged, while every other deletion projection becomes
a matched replacement pair, up to a common permutation. Thus the
property persists. Discarding a pair only makes that entry virtual.
Every two slots occur together in some deletion projection because
$k\ge 2$. The initial types in these projections therefore preserve
all equalities and adjacencies among occupied slots, proving that
this strategy wins.

Conversely, fix a winning strategy with $k+1$ slots. We prove by
induction on $s$ that, at every history ending with exactly $k$
occupied pairs, the two tuples have equal round-$s$ colors in any
common ordering of those slots. The case $s=0$ is partial
isomorphism. For the induction step, fill the spare slot. The
strategy announces a bijection before Spoiler chooses the new
vertex. After any such choice, Spoiler may discard any one of the
old slots. The induction hypothesis at these resulting histories
matches all $k$ replacement colors for that same vertex bijection.
Hence the entire joint vector is matched, and
Equation~\eqref{eq:wl-update} gives the next-round equality.

Fill $k$ slots from empty in a fixed order. This procedure induces a
bijection from $V(G)^k$ to $V(H)^k$, even when the strategy depends
on the history. To recover the preimage of a target tuple, invert
the first announced vertex bijection, thereby recover the first
source choice and history, and continue successively. Since each
resulting tuple pair has equal colors at every round, the color
histograms agree.

For the final assertion, use induction on the ordinary
color-refinement round at all one-pair histories. From such a
history fill another slot. The announced bijection preserves
adjacency to the retained pair for every possible choice. After
discarding the old pair, the induction hypothesis matches the
previous-round colors of the new vertices. Thus the same bijection
matches the neighbor-color multisets of the retained vertices.
Their next colors agree. The first announced bijection, from the
empty history, then proves equality of all vertex-color histograms.
\end{proof}

\begin{lemma}
\label{lem:monotonicity}
For positive integers $j\le k$, the relation $G\equiv_k H$ implies
$G\equiv_j H$. Consequently, a graph has Weisfeiler--Leman dimension
at most $k$ if and only if $k$-WL identifies it.
\end{lemma}

\begin{proof}
Equal histograms imply equal graph orders, since their total sizes
are $|V(G)|^k$ and $|V(H)|^k$. Empty graphs cause no difficulty, so
assume that this common order is positive. A winning strategy can
be restricted to a smaller number of slots. For $k\ge j\ge 2$,
apply Lemma~\ref{lem:wl-game} in both directions. For $k\ge 2$ and
$j=1$, use its final assertion; the remaining case is immediate.
If some $j\le k$ identifies $G$, then $k$-equivalence to $G$
implies $j$-equivalence and hence isomorphism. The converse follows
from the definition of dimension.
\end{proof}

The next elementary observation lets us recover vertex colors that
will be encoded by degrees.

\begin{lemma}
\label{lem:degree}
In every history following a winning strategy with at least two pair
slots, matched vertices have equal degrees. Every vertex bijection
announced by that strategy therefore preserves degrees for all its
possible choices.
\end{lemma}

\begin{proof}
Retain only a suspect pair $(v,w)$ and fill another slot. To preserve
adjacency for every vertex choice, the announced bijection must map
$N_G(v)$ onto $N_H(w)$, so their cardinalities agree. This applies to
every candidate pair created by any announced bijection, proving
the second assertion as well.
\end{proof}

\subsection{An exponential-time decision procedure}

\begin{proposition}
\label{prop:upper}
For a nonempty graph given by its binary adjacency matrix and a
positive integer $k$ in binary, deciding whether its
Weisfeiler--Leman dimension is at most $k$ belongs to EXPTIME.
\end{proposition}

\begin{proof}
Malformed encodings, including matrices outside the specified graph
domain and nonpositive parameters, can be rejected in polynomial
time. Let $n\ge 1$ be the order of a valid input graph $G$.

First compare the binary parameter with $n$. If $k\ge n$, accept.
For $n\ge 2$, a $k$-tuple listing all vertices, with repetitions if
needed, has an initial equality and adjacency type that determines
an isomorphism to any matching tuple in a graph of order $n$.
Different orders are distinguished by histogram sizes. Order one
is immediate under ordinary color refinement. This comparison is
performed before constructing a tuple table, regardless of how
large the integer encoded by $k$ may be.

It remains to consider $1\le k<n$. Enumerate all labeled simple
graphs $H$ on $n$ vertices, and reject exactly when one satisfies
$G\equiv_k H$ but $G\not\cong H$. This is the desired test by
Lemma~\ref{lem:monotonicity}. There are
$2^{\binom n2}$ candidates, and testing all $n!$ vertex permutations
per candidate suffices for isomorphism.

For equivalence, maintain integer identifiers for the common color
classes of the two tuple tables. There are $2n^k$ entries and at
most $2n^k$ refinement rounds before stabilization. Each signature
contains the old identifier and at most $n$ vectors of length $k$;
sorting and renaming these signatures costs a fixed polynomial in
$n^k,n,k$. Recursive color expressions need not be expanded.
For $k=1$ use the neighbor signatures instead. Agreement of the
histograms through stabilization is sufficient, since later rounds
only rename stable classes. Finally, $n^k\le n^n$, so enumeration,
refinement, and permutation testing together take time exponential
in a polynomial of the explicit input length.
\end{proof}
